\documentclass{beamer}
\usepackage{graphicx} % Required for inserting images
\usepackage{float}
\usepackage[T1]{fontenc}
\title{Analysing the Robustness of Deep Neural Networks}
\author{Divij Khaitan}
\date{\today}
\usepackage{amsmath}
\usepackage{amssymb}
\usepackage{subcaption}
\usepackage[backend=biber]{biblatex}
\addbibresource{references.bib}

\DeclareMathOperator{\diag}{diag}
\DeclareMathOperator{\softmax}{softmax}

\begin{document}
    \begin{frame}
        \maketitle
    \end{frame}
    
    % \begin{frame}{Background: Overfitting}
    % \includegraphics[width=0.6\textwidth]{overfitting_example.png}
    % Image credits: Wikipedia
    % \end{frame}
    
    \begin{frame}{Background: Overfitting}
        \begin{figure}[H]
            \centering
            \begin{minipage}{0.3\textwidth}
                \centering
                \includegraphics[width=\linewidth]{underfit.png}
                \subcaption{Underfit}
                % \label{fig:image1}
            \end{minipage}
            \hfill
            \begin{minipage}{0.3\textwidth}
                \centering
                \includegraphics[width=\linewidth]{best_fit.png}
                \subcaption{Best Fit}
                % \label{fig:image2}
            \end{minipage}
            \begin{minipage}{0.3\textwidth}
                \centering
                \includegraphics[width=\linewidth]{overfit.png}
                \subcaption{Overfit}
                % \label{fig:image1}
            \end{minipage}
            
            \caption{Examples of overfitting data from a third degree polynomial}
            \label{fig:overfitting}
        \end{figure}
    \end{frame}

    % \begin{frame}{Overfitting}
    %     \begin{itemize}
    \setlength\itemsep{1em}
    %         \item 
    %             Comes from mathematical modelling
    %         \item 
    %             Very complex models capture spurious relationships
    %         \item 
    %             Rule of thumb $\#(\text{parameters}) > \#(\text{datapoints})$
    %         \item 
    %             Analogous to interpolating a polynomial of degree higher than $n$ with $n$ datapoints
    %     \end{itemize}
    % \end{frame}

    \begin{frame}{Modern ML}
        \begin{itemize}
            \setlength\itemsep{1em}
            \item 
                Alexnet : \#(parameters) $\approx \; 4 \times 10^5$
            \item 
                Achieves about $60\%$ accuracy on Imagenet 1K with $\approx$ 1.3 million images
            \item 
                ResNet-50 ($2.5 \times 10^7$) parameters, $76\%$ accuracy
            \item 
                Is this network overfitting?
        \end{itemize}
    \end{frame}

    \begin{frame}{Rethinking Generalisation}
        \begin{itemize}
            \setlength\itemsep{1em}
            \item 
                Fitting networks to random labels eventually reaches zero training error, even with random pixels (\cite{zhang2017understandingdeeplearningrequires}, \cite{arpit2017closerlookmemorizationdeep})
            \item 
                Models have enough capacity to shatter the training set, yet learn patterns when available
            \item 
                Conclusion: Resistance to overfitting or bias towards simple hypotheses
        \end{itemize}
    \end{frame}

    \begin{frame}{Manifold Hypothesis}

        \includegraphics[width=\textwidth]{manifold.png}
        
        \begin{itemize}
            \setlength\itemsep{1em}
            \item 
                Pre-dates deep learning
            \item 
                Data consists of substructures with much smaller dimension than the ambient space
            \item 
                Valuable as a heuristic, basis of most dimensionality reduction algorithms
        \end{itemize}
    \end{frame}
    
    \begin{frame}{Existing Work on Activation Patterns}
        \begin{itemize}
            \setlength\itemsep{1em}
            \item 
                Interpretability: \cite{stano2020explainingpredictionsdeepneural} fits a gaussian mixture, binarises and uses the hamming distance to support predictions, \cite{bauerle2022neuralactivationpatternsnaps} is a tool to visualise activation patterns as they relate to features
            \item 
                OOD detection: Using specific patterns of activation and associating them with each class, \cite{lee2018simpleunifiedframeworkdetecting} fits a multivariate gaussian to each class in the embedding space and the mahanlanobis distance from an activation pattern to support decisions
            \item 
                Idea of fitting distributions to activations for uncertainty quantification is common
        \end{itemize}
    \end{frame}

    \begin{frame}{Our Contributions}
        \begin{itemize}
            \setlength\itemsep{1em}
            \item 
                Analyse the activations at a finer granularity
            \item 
                Introduce a notion of significance into neuron activations
            \item 
                Distributions over classes provide better characterisation of OOD data than manifold interpretation
            \item 
                Measure for robustness 
        \end{itemize}
    \end{frame}

    \begin{frame}{Neuron-Weight Interactions}
        Every product of a neuron from the previous layer with a weight from the next layer. 
        Given an activation vector $x \in \mathbb{R}^{n_{in} \times 1}$ and a weight matrix $W \in \mathbb{R}^{n_{out} \times n_{in}}$ 
        \[
            \begin{bmatrix}
                w_{1,1} & \dots & w_{1,n_{in}} \\
                \vdots & \ddots & \vdots \\
                w_{n_{out}, 1} & \cdots & w_{n_{out}, n_{in}} \\
            \end{bmatrix}
            \begin{bmatrix}
                x_1 \\
                \vdots \\
                x_{n_{in}} \\
            \end{bmatrix}
            = 
            \begin{bmatrix}
                \sum_{i=1}^{n_in} w_{1,i}x_i\\
                \vdots \\
                \sum_{i=1}^{n_in} w_{n_{out},i}x_i\\
            \end{bmatrix}
        \]
    \end{frame}

    \begin{frame}{Neuron-Weight Interactions}
        To study every term in each summation individually, a vector of 1s needs to be factored out.
        \[
            W\diag(x_1, \dots, x_n)    
            = 
            \begin{bmatrix}
                w_{1, 1}x_1 & \dots & w_{1, n_{in}}x_{n_{in}} \\
                \vdots & \ddots & \vdots \\
                w_{n_{out}, 1}x_1 & \dots & w_{n_{out}, n_{in}}x_{n_{in}}
            \end{bmatrix}
        \]
    \end{frame}
    
    \begin{frame}{Significance of Neuron Weight Interactions}
        An interaction is considered significant if it is large relative to the magnitude of the neuron value, corrected for the input dimension. Formally 
        \[S_{ij} = 1 \text{ if } |x_{ij}| > \frac{|\sum_j x_{ij}|}{n_{in}} \]

    \end{frame}

    \begin{frame}{Activation Distributions}
        \begin{itemize}
            \setlength\itemsep{1em}
            \item 
                Modelled binary on-off values for each interaction in the linear layers as independent bernoulli random variables
            \item 
                Measured KL divergence and entropy for distributions
            \item 
                Can be seen as analogues to the cluster radius and distances
            \item 
                Our current methodology over-estimates the separation between classes
        \end{itemize}
    \end{frame}

    \begin{frame}{Sparsity of Neuron Activations}
        \begin{itemize}
            \setlength\itemsep{1em}
            \item 
                Practically, pruning techniques can drastically decrease the size of neural networks
            \item 
                Indicates that a small part of the network is important
            \item 
                Lottery Ticket Hypothesis and related work focuses on identiying sparse structures early on in training
            \item 
                Activations should be sparse
        \end{itemize}
    \end{frame}


    \begin{frame}{KL Divergence and Entropy}
        \begin{itemize}
            \setlength\itemsep{1em}
            \item 
                Entropy of a distribution is a measure of the amount of information encoded in a random variable
                \[H(X) = -\sum_x p(x)\log(p(x))\]
            \item 
                KL divergence is the relative entropy between 2 variables, i.e. the amount of additional information needed to transform an encoding of P to one of Q
                \[KL(P||Q) = \sum_x P(x)\log(\frac{P(x)}{Q(x)})\]
            \item 
                KL divergence isn't a metric, satisfies nonnegavity but not symmetry or triangle inequality
        \end{itemize}
    \end{frame}

    \begin{frame}{Histogram Distance}
        \includegraphics[width=\textwidth]{histogram_distance.png}
    \end{frame}

    % \begin{frame}{Detecting Adversarial Attacks}
    %     \begin{itemize}
    % \setlength\itemsep{1em}
    %         \item 
    %             Adversarial attacks are small changes to an image that alter it's predicted class, often with a high output probability
    %         \item 
    %             We computed the probabilities of the observed activation pattern under the 
    %     \end{itemize}
    % \end{frame}

    \begin{frame}{Estimating Dimensions of a Manifold}
        \begin{itemize}
            \setlength\itemsep{1em}
            \item 
                Extracted the output of the final hidden layer (embedding), and computed the dimensionality of the manifolds for each class
            \item 
                Final layer activations must be linearly separable 
            \item 
                Computed using principal components upto 95\% explained variance
            \item 
                Measured pairwise entangled classes by computing the rank of the matrix created by the union of PCs
        \end{itemize}
    \end{frame}

    \begin{frame}{Overview of Framework}
        \begin{itemize}
            \item 
                Distributions over neuron-weight interactions: Increase in the separation between classes $\implies$ developing a distinct idea for each class
            \item 
                Sparse Activations: Small number of frequent firings, large number of random firings
            \item 
                OOD Data: Natural probabilistic definition of out-of-distribution data rather than geometric definition as with manifolds
        \end{itemize}
    \end{frame}

    \begin{frame}{Models and Datasets}
        \begin{itemize}
            \setlength\itemsep{1em}
            \item 
                The tests were run on AlexNet, with the pretrained version trained on ImageNet coming from PyTorch
            \item 
                Could not run for higher layers for other networks because of storage restrictions
            \item 
                Compared behaviour of network on OOD dataset imagenet-r \cite{hendrycks2021facesrobustnesscriticalanalysis}
            \item 
                Also reproduced Modified AlexNet experiment on CIFAR10 from \cite{zhang2017understandingdeeplearningrequires}, for random and normal labels
            \item 
                Have 4 versions of modified AlexNet, best model (smallest validation error), overtrained model (run for 250 epochs on clean labels), random labels (0 error on CIFAR10) and untrained (initialised network with no training)
        \end{itemize}
    \end{frame}

    \begin{frame}{Results: Sparsity - CIFAR10}
        \begin{figure}[H]
            \centering
            \begin{minipage}{0.45\textwidth}
                \centering
                \includegraphics[width=\linewidth]{new/small_alexnet_untrained_cifar10/activation_proportions/0.png}
                \subcaption{Untrained}
                % \label{fig:image1}
            \end{minipage}
            \hfill
            \begin{minipage}{0.45\textwidth}
                \centering
                \includegraphics[width=\linewidth]{new/small_alexnet_cifar10_random_labels/activation_proportions/0.png}
                \subcaption{Random Labels}
                % \label{fig:image2}
            \end{minipage}
            \begin{minipage}{0.45\textwidth}
                \centering
                \includegraphics[width=\linewidth]{new/small_alexnet_cifar10_250_epoch/activation_proportions/0.png}
                \subcaption{Overtrained}
                % \label{fig:image1}
            \end{minipage}
            \hfill
            \begin{minipage}{0.45\textwidth}
                \centering
                \includegraphics[width=\linewidth]{new/small_alexnet_cifar10_best/activation_proportions/0.png}
                \subcaption{Best}
                % \label{fig:image2}
            \end{minipage}
            \hfill
            
            \caption{Activation Proportions from the class Airplane on the last and second last layers on different versions of the Modified AlexNet}
            \label{fig:class_airplane}
        \end{figure}
    \end{frame}

    \begin{frame}{Results: Sparsity - CIFAR10}
        \begin{itemize}
            \setlength\itemsep{1em}
            \item 
                Training encourages sparsity 
            \item 
                Proportions of neurons that are always active drops 
            \item     
                Random labels has an extremely sparse penultimate layer, best is still sparse but denser than random 
            \item 
                Overtrained network a mid-point in between random labels and best
        \end{itemize}
    \end{frame}
    
    \begin{frame}{Results: Sparsity - Imagenet and imagenet-r}
        \begin{figure}[H]
            \centering
            \begin{minipage}{0.45\textwidth}
                \centering
                \includegraphics[width=\linewidth]{new/alexnet_untrained_imgr_out/activation_proportions/n01677366.png}
                \subcaption{Untrained Imagenet}
                % \label{fig:image1}
            \end{minipage}
            \hfill
            \begin{minipage}{0.45\textwidth}
                \centering
                \includegraphics[width=\linewidth]{new/alexnet_imgr_out/activation_proportions/n01677366.png}
                \subcaption{Trained Imagenet}
                % \label{fig:image2}
            \end{minipage}
            \begin{minipage}{0.45\textwidth}
                \centering
                \includegraphics[width=\linewidth]{new/alexnet_untrained_r/activation_proportions/n01677366.png}
                \subcaption{Untrained imagenet-r}
                % \label{fig:image1}
            \end{minipage}
            \hfill
            \begin{minipage}{0.45\textwidth}
                \centering
                \includegraphics[width=\linewidth]{new/alexnet_r/activation_proportions/n01677366.png}
                \subcaption{Trained imagenet-r}
                % \label{fig:image2}
            \end{minipage}
            \hfill
            
            \caption{Activation Proportions from the class iguana on the last three layers on images from imagenet and imagenet-r on the trained and untrained alexnet models}
            \label{fig:alexnet_sparsity}
        \end{figure}
    \end{frame}

    \begin{frame}{Results: Sparsity - Imagenet}
        \begin{itemize}
            \setlength\itemsep{1em}
            \item 
                Untrained looks the same for ID and OOD
            \item 
                ID set sparser on last two layers 
            \item     
                More stray activations on the last two layers for the OOD set 
        \end{itemize}
    \end{frame}
    
    \begin{frame}{Results: Distances between distributions - CIFAR10}
        \begin{figure}[H]
            \centering
            \begin{minipage}{0.24\textwidth}
                \centering
                \includegraphics[width=\linewidth]{new/small_alexnet_untrained_cifar10/out_matrices/small_alexnet_kl_matrix.png}
                \subcaption{Untrained}
                % \label{fig:image1}
            \end{minipage}
            \hfill
            \begin{minipage}{0.24\textwidth}
                \centering
                \includegraphics[width=\linewidth]{new/small_alexnet_cifar10_random_labels/out_matrices/small_alexnet_kl_matrix.png}
                \subcaption{Random Labels}
                % \label{fig:image2}
            \end{minipage}
            \begin{minipage}{0.24\textwidth}
                \centering
                \includegraphics[width=\linewidth]{new/small_alexnet_cifar10_250_epoch/out_matrices/small_alexnet_kl_matrix.png}
                \subcaption{Overtrained}
                % \label{fig:image1}
            \end{minipage}
            \hfill
            \begin{minipage}{0.24\textwidth}
                \centering
                \includegraphics[width=\linewidth]{new/small_alexnet_cifar10_best/out_matrices/small_alexnet_kl_matrix.png}
                \subcaption{Best}
                % \label{fig:image2}
            \end{minipage}
            \hfill
            
            \caption{Activation Proportions from the class Airplane on the last and second last layers on different versions of the Modified AlexNet}
            \label{fig:small_alexnet_kl_matrix}
        \end{figure}
    \end{frame}
    
    \begin{frame}{Results: Distances between distributions - CIFAR10}
        \begin{figure}[H]
            \centering
            \begin{minipage}{0.24\textwidth}
                \centering
                \includegraphics[width=\linewidth]{new/small_alexnet_untrained_cifar10_second_last/out_matrices/small_alexnet_kl_matrix.png}
                \subcaption{Untrained}
                % \label{fig:image1}
            \end{minipage}
            \hfill
            \begin{minipage}{0.24\textwidth}
                \centering
                \includegraphics[width=\linewidth]{new/small_alexnet_cifar10_random_labels_second_last/out_matrices/small_alexnet_kl_matrix.png}
                \subcaption{Random Labels}
                % \label{fig:image2}
            \end{minipage}
            \begin{minipage}{0.24\textwidth}
                \centering
                \includegraphics[width=\linewidth]{new/small_alexnet_cifar10_250_epoch_second_last/out_matrices/small_alexnet_kl_matrix.png}
                \subcaption{Overtrained}
                % \label{fig:image1}
            \end{minipage}
            \hfill
            \begin{minipage}{0.24\textwidth}
                \centering
                \includegraphics[width=\linewidth]{new/small_alexnet_cifar10_best_second_last/out_matrices/small_alexnet_kl_matrix.png}
                \subcaption{Best}
                % \label{fig:image2}
            \end{minipage}
            \hfill
            
            \caption{Activation Proportions from the class Airplane on the last and second last layers on different versions of the Modified AlexNet}
            \label{fig:small_alexnet_kl_matrix_second_last}
        \end{figure}
    \end{frame}

    \begin{frame}{Results: Changes in Divergence - CIFAR10}
        \begin{figure}[H]
            \centering
            \begin{minipage}{0.22\textwidth}
                \centering
                \includegraphics[width=\linewidth]{new/untrained_random_cifar10.png}
                \subcaption{Untrained and Random Last Layer}
                % \label{fig:image1}
            \end{minipage}
            \hfill
            \begin{minipage}{0.22\textwidth}
                \centering
                \includegraphics[width=\linewidth]{new/best_overtrained_cifar10.png}
                \subcaption{Best and Overtrained Last Layer}
                % \label{fig:image2}
            \end{minipage}
            \hfill
            \begin{minipage}{0.22\textwidth}
                \centering
                \includegraphics[width=\linewidth]{new/untrained_random_cifar10_second_last.png}
                \subcaption{Untrained and Random Second Last Layer}
                % \label{fig:image1}
            \end{minipage}
            \hfill
            \begin{minipage}{0.22\textwidth}
                \centering
                \includegraphics[width=\linewidth]{new/best_overtrained_cifar10_second_last.png}
                \subcaption{Best and Overtrained Second Last Layer}
                % \label{fig:image2}
            \end{minipage}
            \hfill
            \caption{Changes in the KL divergence over the course of training}
            \label{fig:small_alexnet_kl_differences}
        \end{figure}
    \end{frame}
    

    \begin{frame}{Results: Distances Between Distributions}
        \begin{itemize}
            \setlength\itemsep{1em}
            \item 
                Training makes classes more separated
            \item 
                Entropy of distributions decreases in lower layers, increases in upper layers
            \item 
                Overtraining and normal labels look indentical with a lower spread for the overtrained network, suggesting that overtraining makes the distributions more similar
            \item 
                Random labels and untrained have no relative difference in the separation and entropy of distributions. 
        \end{itemize}        
    \end{frame}

    \begin{frame}{Results: Distances between distributions - Imagenet}
        % Pairwise heatmaps for the last layer 
        \begin{figure}[H]
            \centering
            \begin{minipage}{0.30\textwidth}
                \centering
                \includegraphics[width=\linewidth]{new/alexnet_untrained_imgr_out/out_matrices/alexnet_kl_matrix.png}
                \subcaption{Untrained Last}
                % \label{fig:image1}
            \end{minipage}
            \hfill
            \begin{minipage}{0.30\textwidth}
                \centering
                \includegraphics[width=\linewidth]{new/alexnet_untrained_imgr_out_second_last/out_matrices/alexnet_kl_matrix.png}
                \subcaption{Untrained Second Last}
                % \label{fig:image1}
            \end{minipage}
            \hfill
            \begin{minipage}{0.30\textwidth}
                \centering
                \includegraphics[width=\linewidth]{new/alexnet_untrained_imgr_out_third_last/out_matrices/alexnet_kl_matrix.png}
                \subcaption{Untrained Third Last}
                % \label{fig:image1}
            \end{minipage}
            \hfill
            \begin{minipage}{0.30\textwidth}
                \centering
                \includegraphics[width=\linewidth]{new/alexnet_imgr_out/out_matrices/alexnet_kl_matrix.png}
                \subcaption{Trained Last}
                % \label{fig:image2}
            \end{minipage}
            \hfill
            \begin{minipage}{0.30\textwidth}
                \centering
                \includegraphics[width=\linewidth]{new/alexnet_imgr_out_second_last/out_matrices/alexnet_kl_matrix.png}
                \subcaption{Trained Second Last}
                % \label{fig:image1}
            \end{minipage}
            \hfill
            \begin{minipage}{0.30\textwidth}
                \centering
                \includegraphics[width=\linewidth]{new/alexnet_imgr_out_third_last/out_matrices/alexnet_kl_matrix.png}
                \subcaption{Trained Third Last}
                % \label{fig:image1}
            \end{minipage}
            \hfill
            
            \caption{KL Divergence Heatmaps on imagenet for Alexnet}
            \label{fig:alexnet_imgr_heatmaps}
        \end{figure}
    \end{frame}
    
    \begin{frame}{Results: Changes in Divergence - Imagenet}
        % Pairwise heatmaps for the last layer 
        \begin{figure}[H]
            \centering
            \begin{minipage}{0.30\textwidth}
                \centering
                \includegraphics[width=\linewidth]{new/untrained_trained_imgr.png}
                \subcaption{Last Layer}
                % \label{fig:image1}
            \end{minipage}
            \hfill
            \begin{minipage}{0.30\textwidth}
                \centering
                \includegraphics[width=\linewidth]{new/untrained_trained_imgr_second_last.png}
                \subcaption{Second Last Layer}
                % \label{fig:image1}
            \end{minipage}
            \hfill
            \begin{minipage}{0.30\textwidth}
                \centering
                \includegraphics[width=\linewidth]{new/untrained_trained_imgr_third_last.png}
                \subcaption{Third Last Layer}
                % \label{fig:image1}
            \end{minipage}
            \caption{Changes in KL divergence on imagenet for Alexnet}
            \label{fig:change_alexnet_imgr_heatmaps}
        \end{figure}
    \end{frame}
    

    \begin{frame}{Results: Distances between distributions - imagenet-r}
        % Pairwise heatmaps for the last layer 
        \begin{figure}[H]
            \centering
            \begin{minipage}{0.30\textwidth}
                \centering
                \includegraphics[width=\linewidth]{new/alexnet_untrained_r/out_matrices/alexnet_kl_matrix.png}
                \subcaption{Untrained Last}
                % \label{fig:image1}
            \end{minipage}
            \hfill
            \begin{minipage}{0.30\textwidth}
                \centering
                \includegraphics[width=\linewidth]{new/alexnet_untrained_r_second_last/out_matrices/alexnet_kl_matrix.png}
                \subcaption{Untrained Second Last}
                % \label{fig:image1}
            \end{minipage}
            \hfill
            \begin{minipage}{0.30\textwidth}
                \centering
                \includegraphics[width=\linewidth]{new/alexnet_untrained_r_third_last/out_matrices/alexnet_kl_matrix.png}
                \subcaption{Untrained Third Last}
                % \label{fig:image1}
            \end{minipage}
            \hfill
            \begin{minipage}{0.30\textwidth}
                \centering
                \includegraphics[width=\linewidth]{new/alexnet_r/out_matrices/alexnet_kl_matrix.png}
                \subcaption{Trained Last}
                % \label{fig:image2}
            \end{minipage}
            \hfill
            \begin{minipage}{0.30\textwidth}
                \centering
                \includegraphics[width=\linewidth]{new/alexnet_r_second_last/out_matrices/alexnet_kl_matrix.png}
                \subcaption{Trained Second Last}
                % \label{fig:image1}
            \end{minipage}
            \hfill
            \begin{minipage}{0.30\textwidth}
                \centering
                \includegraphics[width=\linewidth]{new/alexnet_r_third_last/out_matrices/alexnet_kl_matrix.png}
                \subcaption{Trained Third Last}
                % \label{fig:image1}
            \end{minipage}
            \hfill
            
            \caption{KL Divergence Heatmaps on imagenet-r for Alexnet}
            \label{fig:alexnet_r_heatmaps}
        \end{figure}
    \end{frame}
    
    \begin{frame}{Results: Changes in Divergence - imagenet-r}
        % Pairwise heatmaps for the last layer 
        \begin{figure}[H]
            \centering
            \begin{minipage}{0.30\textwidth}
                \centering
                \includegraphics[width=\linewidth]{new/untrained_trained_r.png}
                \subcaption{Last Layer}
                % \label{fig:image1}
            \end{minipage}
            \hfill
            \begin{minipage}{0.30\textwidth}
                \centering
                \includegraphics[width=\linewidth]{new/untrained_trained_r_second_last.png}
                \subcaption{Second Last Layer}
                % \label{fig:image1}
            \end{minipage}
            \hfill
            \begin{minipage}{0.30\textwidth}
                \centering
                \includegraphics[width=\linewidth]{new/untrained_trained_r_third_last.png}
                \subcaption{Third Last Layer}
                % \label{fig:image1}
            \end{minipage}
            \caption{Changes in KL divergence on imagenet-r for Alexnet}
            \label{fig:change_alexnet_r_heatmaps}
        \end{figure}
    \end{frame}
    

    \begin{frame}{Results: Distances Between Distributions - Imagenet}
        \begin{itemize}
            \setlength\itemsep{1em}
            \item 
               Entropy of OOD class distributions sees no significant drop after training
            \item 
                The OOD distributions increase in separation, but the magnitude of the change is much smaller
        \end{itemize}        
    \end{frame}

    % \begin{frame}{Results: Paiwise Distances between distributions on Imagenet-R and Imagenet}
    %     Pairwise heatmaps for the last layer 
        
    % \end{frame}
    
    \begin{frame}{Results: Distances between distribution histograms - CIFAR10}
        \begin{figure}[H]
            \centering
            \begin{minipage}{0.24\textwidth}
                \centering
                \includegraphics[width=\linewidth]{new/small_alexnet_untrained_cifar10/out_matrices/small_alexnet_intersection_matrix.png}
                \subcaption{Untrained}
                % \label{fig:image1}
            \end{minipage}
            \hfill
            \begin{minipage}{0.24\textwidth}
                \centering
                \includegraphics[width=\linewidth]{new/small_alexnet_cifar10_random_labels/out_matrices/small_alexnet_intersection_matrix.png}
                \subcaption{Random Labels}
                % \label{fig:image2}
            \end{minipage}
            \begin{minipage}{0.24\textwidth}
                \centering
                \includegraphics[width=\linewidth]{new/small_alexnet_cifar10_250_epoch/out_matrices/small_alexnet_intersection_matrix.png}
                \subcaption{Overtrained}
                % \label{fig:image1}
            \end{minipage}
            \hfill
            \begin{minipage}{0.24\textwidth}
                \centering
                \includegraphics[width=\linewidth]{new/small_alexnet_cifar10_best/out_matrices/small_alexnet_intersection_matrix.png}
                \subcaption{Best}
                % \label{fig:image2}
            \end{minipage}
            \hfill
            
            \caption{Activation Proportions from the class Airplane on the last and second last layers on different versions of the Modified AlexNet}
            \label{fig:small_alexnet_intersection_matrix}
        \end{figure}
    \end{frame}
    
    \begin{frame}{Results: Distances between distribution histograms - CIFAR10}
        \begin{figure}[H]
            \centering
            \begin{minipage}{0.24\textwidth}
                \centering
                \includegraphics[width=\linewidth]{new/small_alexnet_untrained_cifar10_second_last/out_matrices/small_alexnet_intersection_matrix.png}
                \subcaption{Untrained}
                % \label{fig:image1}
            \end{minipage}
            \hfill
            \begin{minipage}{0.24\textwidth}
                \centering
                \includegraphics[width=\linewidth]{new/small_alexnet_cifar10_random_labels_second_last/out_matrices/small_alexnet_intersection_matrix.png}
                \subcaption{Random Labels}
                % \label{fig:image2}
            \end{minipage}
            \begin{minipage}{0.24\textwidth}
                \centering
                \includegraphics[width=\linewidth]{new/small_alexnet_cifar10_250_epoch_second_last/out_matrices/small_alexnet_intersection_matrix.png}
                \subcaption{Overtrained}
                % \label{fig:image1}
            \end{minipage}
            \hfill
            \begin{minipage}{0.24\textwidth}
                \centering
                \includegraphics[width=\linewidth]{new/small_alexnet_cifar10_best_second_last/out_matrices/small_alexnet_intersection_matrix.png}
                \subcaption{Best}
                % \label{fig:image2}
            \end{minipage}
            \hfill
            
            \caption{Activation Proportions from the class Airplane on the last and second last layers on different versions of the Modified AlexNet}
            \label{fig:small_alexnet_intersection_matrix_second_last}
        \end{figure}
    \end{frame}

    \begin{frame}{Results: Distances between distribution histograms - Imagenet}
        % Pairwise heatmaps for the last layer 
        \begin{figure}[H]
            \centering
            \begin{minipage}{0.30\textwidth}
                \centering
                \includegraphics[width=\linewidth]{new/alexnet_untrained_imgr_out/out_matrices/alexnet_intersection_matrix.png}
                \subcaption{Untrained Last}
                % \label{fig:image1}
            \end{minipage}
            \hfill
            \begin{minipage}{0.30\textwidth}
                \centering
                \includegraphics[width=\linewidth]{new/alexnet_untrained_imgr_out_second_last/out_matrices/alexnet_intersection_matrix.png}
                \subcaption{Untrained Second Last}
                % \label{fig:image1}
            \end{minipage}
            \hfill
            \begin{minipage}{0.30\textwidth}
                \centering
                \includegraphics[width=\linewidth]{new/alexnet_untrained_imgr_out_third_last/out_matrices/alexnet_intersection_matrix.png}
                \subcaption{Untrained Third Last}
                % \label{fig:image1}
            \end{minipage}
            \hfill
            \begin{minipage}{0.30\textwidth}
                \centering
                \includegraphics[width=\linewidth]{new/alexnet_imgr_out/out_matrices/alexnet_intersection_matrix.png}
                \subcaption{Trained Last}
                % \label{fig:image2}
            \end{minipage}
            \hfill
            \begin{minipage}{0.30\textwidth}
                \centering
                \includegraphics[width=\linewidth]{new/alexnet_imgr_out_second_last/out_matrices/alexnet_intersection_matrix.png}
                \subcaption{Trained Second Last}
                % \label{fig:image1}
            \end{minipage}
            \hfill
            \begin{minipage}{0.30\textwidth}
                \centering
                \includegraphics[width=\linewidth]{new/alexnet_imgr_out_third_last/out_matrices/alexnet_intersection_matrix.png}
                \subcaption{Trained Third Last}
                % \label{fig:image1}
            \end{minipage}
            \hfill
            
            \caption{KL Divergence Heatmaps on imagenet for Alexnet}
            \label{fig:alexnet_imgr_histogram_distance}
        \end{figure}
    \end{frame}
    
    \begin{frame}{Results: Distances between distribution histograms - imagenet-r}
        % Pairwise heatmaps for the last layer 
        \begin{figure}[H]
            \centering
            \begin{minipage}{0.30\textwidth}
                \centering
                \includegraphics[width=\linewidth]{new/alexnet_untrained_r/out_matrices/alexnet_intersection_matrix.png}
                \subcaption{Untrained Last}
                % \label{fig:image1}
            \end{minipage}
            \hfill
            \begin{minipage}{0.30\textwidth}
                \centering
                \includegraphics[width=\linewidth]{new/alexnet_untrained_r_second_last/out_matrices/alexnet_intersection_matrix.png}
                \subcaption{Untrained Second Last}
                % \label{fig:image1}
            \end{minipage}
            \hfill
            \begin{minipage}{0.30\textwidth}
                \centering
                \includegraphics[width=\linewidth]{new/alexnet_untrained_r_third_last/out_matrices/alexnet_intersection_matrix.png}
                \subcaption{Untrained Third Last}
                % \label{fig:image1}
            \end{minipage}
            \hfill
            \begin{minipage}{0.30\textwidth}
                \centering
                \includegraphics[width=\linewidth]{new/alexnet_r/out_matrices/alexnet_intersection_matrix.png}
                \subcaption{Trained Last}
                % \label{fig:image2}
            \end{minipage}
            \hfill
            \begin{minipage}{0.30\textwidth}
                \centering
                \includegraphics[width=\linewidth]{new/alexnet_r_second_last/out_matrices/alexnet_intersection_matrix.png}
                \subcaption{Trained Second Last}
                % \label{fig:image1}
            \end{minipage}
            \hfill
            \begin{minipage}{0.30\textwidth}
                \centering
                \includegraphics[width=\linewidth]{new/alexnet_r_third_last/out_matrices/alexnet_intersection_matrix.png}
                \subcaption{Trained Third Last}
                % \label{fig:image1}
            \end{minipage}
            \hfill
            
            \caption{KL Divergence Heatmaps on imagenet-r for Alexnet}
            \label{fig:alexnet_r_histogram_distance}
        \end{figure}
    \end{frame}

    \begin{frame}{Results: Dimensions of Manifolds on CIFAR10}
        \begin{figure}[H]
            \centering
            \begin{minipage}{0.30\textwidth}
                \centering
                \includegraphics[width=\linewidth]{new/small_alexnet_cifar10_random_labels/probability_results/pca_overlap_heatmap.png}
                \subcaption{Random Labels}
                % \label{fig:image1}
            \end{minipage}
            \hfill
            \begin{minipage}{0.30\textwidth}
                \centering
                \includegraphics[width=\linewidth]{new/small_alexnet_cifar10_250_epoch/probability_results/pca_overlap_heatmap.png}
                \subcaption{Overtrained}
                % \label{fig:image1}
            \end{minipage}
            \hfill
            \begin{minipage}{0.30\textwidth}
                \centering
                \includegraphics[width=\linewidth]{new/small_alexnet_cifar10_best/probability_results/pca_overlap_heatmap.png}
                \subcaption{Best}
                % \label{fig:image1}
            \end{minipage}
            \hfill
            \caption{Manifold Dimension Overlap Heatmaps}
            \label{fig:pca_overlap_heatmaps}
        \end{figure}
    \end{frame}

    \begin{frame}{Dimensionality of Manifolds}
        \begin{itemize}
            \setlength\itemsep{1em}
            \item 
                Significant overlap on the random label setting
            \item 
                Much smaller dimensionality on actual data
            \item 
                Overtrained model has higher dimensional manifolds for every single class
        \end{itemize}    
    \end{frame}
    
    \begin{frame}{Conclusions}
        \begin{itemize}
            \setlength\itemsep{1em}
            \item 
                Significant activations are sparse
            \item 
                Framework registers different behaviour on in and out of distribution datasets
            \item 
                Behaviour also different for networks that memorise and overtrain compared to a best-fit network
        \end{itemize}
    \end{frame}
    
    \begin{frame}{Future Work}
        \begin{itemize}
            \setlength\itemsep{1em}
            \item 
                Sample complexity of attaining a KL Divergence of $\epsilon$
            \item 
                Measuring data regularity/complexity based on the KL divergence
            \item 
                Tightening the esimate by making a modification/finding an under-estimator
        \end{itemize}
    \end{frame}

    \begin{frame}[allowframebreaks]{References}
        \printbibliography
    \end{frame}
        
\end{document}